\begin{lemma}[Near-diagonal size, gradient and H\"ormander bounds for $K_{\gamma,j}$]
Fix $n \geq 1$ and $j \in \{1,\ldots,n\}$, let $\lambda$ be Lebesgue measure on
$\mathbb{R}^n$, and let $K_{\gamma,j}$ be the normalized Gaussian Riesz kernel
of \noderef{GaussRieszKernel}, so that
$K_{\gamma,j}(x,y) = \sqrt2\,\pi^{-(n+1)/2}\int_0^1\Phi_{x,y}(r)\,dr$ with
\[
  \Phi_{x,y}(r)
    = \frac{y_j - r x_j}{(1-r^2)^{n/2+1}\sqrt{-\log r}}
      \exp\Bigl(-\frac{|rx-y|^2}{1-r^2}\Bigr),
  \qquad r \in (0,1) .
\]
Let $R > 0$.  Then there is a constant $C > 0$, depending only on $n$, $j$ and
$R$, such that:
\begin{enumerate}
\item for all $x,y \in \mathbb{R}^n$ with $|x| \leq R+1$ and
$0 < |x-y| \leq 2$,
  \begin{enumerate}
  \item $\Phi_{x,y}$ is Lebesgue integrable on $(0,1)$;
  \item $\displaystyle |K_{\gamma,j}(x,y)| \leq \frac{C}{|x-y|^{\,n}}$; and
  \item the function $w \mapsto K_{\gamma,j}(w,y)$ is (Fr\'echet)
    differentiable at $x$, and its derivative $g$ there --- a continuous linear
    functional on $\mathbb{R}^n$ --- satisfies
    $\displaystyle \|g\| \leq \frac{C}{|x-y|^{\,n+1}}$;
  \end{enumerate}
\item (H\"ormander estimate) for all $x,y,z \in \mathbb{R}^n$ with
$|x| \leq R$, $0 < |x-y| \leq 1$ and $|x-z| \leq \tfrac12|x-y|$,
\[
  \bigl|K_{\gamma,j}(x,y) - K_{\gamma,j}(z,y)\bigr|
    \leq \frac{C\,|x-z|}{|x-y|^{\,n+1}} .
\]
\end{enumerate}
Since $|x-y| > 0$ throughout, the quotients $C/|x-y|^n$ and $C/|x-y|^{n+1}$ are
the usual negative powers $C|x-y|^{-n}$ and $C|x-y|^{-n-1}$.

\emph{Why (1)(a) and the differentiability in (1)(c) are part of the statement,
and not merely of the proof.}  $K_{\gamma,j}(x,y)$ is a Bochner integral over
$(0,1)$ whose value is $0$ by convention when the integrand is not integrable,
and $0$ satisfies every upper bound; so a bare size bound $|K_{\gamma,j}| \leq
C|x-y|^{-n}$ would be provable with no information about the kernel.  Clause
(1)(a) removes that, and the reference paper's own display of
$k_j^{\mathrm{raw}}$ carries the qualifier $x \neq y$, so the integrability
assertion records what the paper intends rather than strengthening it.
Symmetrically, a gradient bound phrased as a bound on a \emph{derivative
operator} chosen by convention where the function is not differentiable would
be vacuous wherever differentiability fails, and would not support a mean value
argument.  Clause (1)(c) therefore asserts genuine differentiability and bounds
the actual derivative.

\emph{Why the H\"ormander clause carries its own hypotheses.}  It is not a
formal consequence of (1)(c): the mean value inequality needs a \emph{uniform}
bound on the derivative along the whole segment from $x$ to $z$, whereas (1)(c)
gives a bound varying with the base point.  The hypothesis
$|x-z| \leq \tfrac12|x-y|$ is exactly what keeps that segment at distance at
least $\tfrac12|x-y|$ from $y$, so that the varying bound is uniformly
comparable to $|x-y|^{-n-1}$ there; the hypotheses $|x| \leq R$ and
$|x-y| \leq 1$ keep the whole segment inside the range of clause (1).

\emph{The constant depends on $R$, necessarily.}  After the splitting of the
exponent used below, the kernel carries a factor $e^{|x|^2-|y|^2}$, which is
bounded only on bounded sets of $x$; on the range considered it is at most
$e^{(R+3)^2}$.  Neither bound is stated with a dimension-only constant, and no
uniformity in $R$ is claimed.
\end{lemma}
\begin{proof}
Write $c_n = \sqrt2\,\pi^{-(n+1)/2}$ and $M = R+3 \geq 3$.  Throughout
$\tau = 1 - r$, so that $r \in (0,1)$ corresponds to $\tau \in (0,1)$, and the
substitution $r \mapsto 1-r$ preserves Lebesgue measure on $\mathbb{R}$ (it is
a reflection followed by a translation; shared context).

\emph{Step 0: the splitting identity and the elementary $\tau$-bounds.}
For all $w,y \in \mathbb{R}^n$ and $r \in (0,1)$,
\begin{equation*}
  \frac{|rw-y|^2}{1-r^2} = \bigl(|y|^2-|w|^2\bigr) + \frac{|w-ry|^2}{1-r^2} ,
  \tag{1}
\end{equation*}
since after multiplying by $1-r^2 > 0$ both
$|rw-y|^2 - (1-r^2)(|y|^2-|w|^2)$ and $|w-ry|^2$ expand to
$|w|^2 - 2r\langle w,y\rangle + r^2|y|^2$.  Hence, whenever $|w| \leq M$,
\begin{equation*}
  \exp\Bigl(-\frac{|rw-y|^2}{1-r^2}\Bigr)
    = e^{|w|^2}e^{-|y|^2}\exp\Bigl(-\frac{|w-ry|^2}{1-r^2}\Bigr)
    \leq e^{M^2}\exp\Bigl(-\frac{|w-ry|^2}{1-r^2}\Bigr) , \tag{2}
\end{equation*}
using $e^{-|y|^2} \leq 1$.  Also $1-r^2 = \tau(1+r) \in [\tau,2\tau]$ and
$\log r = \log(1-\tau) \leq -\tau$, so $-\log r \geq \tau$ and
\begin{equation*}
  \frac{1}{(1-r^2)^{n/2+1}\sqrt{-\log r}} \leq \tau^{-n/2-3/2} ,
  \qquad
  \exp\Bigl(-\frac{|w-ry|^2}{1-r^2}\Bigr)
    \leq \exp\Bigl(-\frac{|w-ry|^2}{2\tau}\Bigr) . \tag{3}
\end{equation*}
We record two integral facts.  For $a > 1$ and $b > 0$, the substitution
$\tau = b/u$ (so $d\tau = (b/u^2)du$) gives
\begin{equation*}
  \int_0^\infty \tau^{-a}e^{-b/\tau}\,d\tau
    = b^{1-a}\int_0^\infty u^{a-2}e^{-u}\,du
    = b^{1-a}\,\Gamma(a-1) < \infty , \tag{4}
\end{equation*}
the last integral converging because $a-2 > -1$ and $e^{-u}$ decays.  And for
$0 < a \leq 1$ and real $p \leq q$ with $q > 1$, using $\tau^{-p} \leq
\tau^{-q}$ on $(0,1]$,
\begin{equation*}
  \int_a^1\tau^{-p}\,d\tau \leq \int_a^1\tau^{-q}\,d\tau
    = \frac{a^{1-q}-1}{q-1} \leq \frac{a^{1-q}}{q-1} . \tag{5}
\end{equation*}

\emph{Step 1: pointwise bounds on $\Phi$ and on its $w$-gradient.}
Fix $y$ and let $w$ satisfy $|w| \leq M$ and $|y| \leq M$; write
$d = |w-y|$.  Since $y_j - rw_j = (y_j-w_j) + \tau w_j$ and
$y - rw = (y-w) + \tau w$,
\begin{equation*}
  |y_j - rw_j| \leq d + \tau M, \qquad |rw - y| \leq d + \tau M . \tag{6}
\end{equation*}
Combining $(2)$, $(3)$ and $(6)$,
\begin{equation*}
  |\Phi_{w,y}(r)|
    \leq e^{M^2}\,(d+\tau M)\,\tau^{-n/2-3/2}
      \exp\Bigl(-\frac{|w-ry|^2}{2\tau}\Bigr) . \tag{7}
\end{equation*}
For each fixed $r \in (0,1)$ the map $w \mapsto \Phi_{w,y}(r)$ is smooth on
$\mathbb{R}^n$: it is a product of a polynomial in $w$ with the exponential of a
polynomial in $w$, divided by the positive constants $(1-r^2)^{n/2+1}$ and
$\sqrt{-\log r}$.  Differentiating, with $e_j$ the $j$th standard basis vector,
\[
  \nabla_w\Phi_{w,y}(r)
    = \frac{\exp\bigl(-|rw-y|^2/(1-r^2)\bigr)}{(1-r^2)^{n/2+1}\sqrt{-\log r}}
      \Bigl[\,-r\,e_j - (y_j - rw_j)\,\frac{2r(rw-y)}{1-r^2}\,\Bigr] ,
\]
because $\nabla_w(y_j-rw_j) = -re_j$ and
$\nabla_w\bigl(-|rw-y|^2/(1-r^2)\bigr) = -2r(rw-y)/(1-r^2)$.  Using
$r \leq 1$, $1-r^2 \geq \tau$, $(2)$, $(3)$ and $(6)$,
\begin{equation*}
  |\nabla_w\Phi_{w,y}(r)|
    \leq e^{M^2}\,\tau^{-n/2-3/2}\Bigl(1 + \frac{2(d+\tau M)^2}{\tau}\Bigr)
      \exp\Bigl(-\frac{|w-ry|^2}{2\tau}\Bigr) . \tag{8}
\end{equation*}

\emph{Step 2: the size bound and clause (1)(a)--(1)(b).}
Let $|x| \leq R+1$ and $0 < d := |x-y| \leq 2$.  Then $|x| \leq M$ and
$|y| \leq |x| + d \leq R+3 = M$, so $(7)$ applies with $w = x$.  Moreover
$x - ry = (x-y) + \tau y$ gives $|x-ry| \geq d - \tau M$.  Put
$\tau' = d/(2M)$; since $d \leq 2$ and $M \geq 3$ we have
$\tau' \leq \tfrac13 < 1$.

For $0 < \tau \leq \tau'$ we have $\tau M \leq d/2$, hence $d + \tau M \leq
\tfrac32 d$ and $|x-ry| \geq d/2$, so $(7)$ gives
$|\Phi_{x,y}| \leq e^{M^2}\tfrac32 d\,\tau^{-n/2-3/2}e^{-d^2/(8\tau)}$, and by
$(4)$ with $a = \tfrac{n+3}{2}$ and $b = d^2/8$,
\[
  \int_{\{0<\tau\leq\tau'\}}|\Phi_{x,y}|\,d\tau
    \leq e^{M^2}\tfrac32 d\,(d^2/8)^{-(n+1)/2}\Gamma\bigl(\tfrac{n+1}{2}\bigr)
    = e^{M^2}\tfrac32\,8^{(n+1)/2}\Gamma\bigl(\tfrac{n+1}{2}\bigr)\,d^{-n} .
\]
For $\tau' < \tau < 1$ we discard the exponential in $(7)$, bounding it by $1$,
and expand the prefactor:
$(d+\tau M)\tau^{-n/2-3/2} = d\,\tau^{-n/2-3/2} + M\,\tau^{-n/2-1/2}$, so that
$|\Phi_{x,y}| \leq e^{M^2}\bigl(d\,\tau^{-n/2-3/2} + M\,\tau^{-n/2-1/2}\bigr)$.
By $(5)$ with $p = q = \tfrac{n+3}{2} > 1$ and with
$p = \tfrac{n+1}{2} \leq q = \tfrac{n+2}{2} > 1$ respectively,
\[
  \int_{\tau'}^1 d\,\tau^{-n/2-3/2}\,d\tau \leq \frac{2d}{n+1}(\tau')^{-(n+1)/2}
    = \frac{2(2M)^{(n+1)/2}}{n+1}\,d^{(1-n)/2} ,
  \qquad
  \int_{\tau'}^1 M\,\tau^{-n/2-1/2}\,d\tau \leq \frac{2M(2M)^{n/2}}{n}\,d^{-n/2} .
\]
Since $0 < d \leq 2$ we have $d^{(1-n)/2} = d^{-n}d^{(n+1)/2} \leq
2^{(n+1)/2}d^{-n}$ and $d^{-n/2} = d^{-n}d^{n/2} \leq 2^{n/2}d^{-n}$.  Adding
the three contributions,
\begin{equation*}
  \int_0^1|\Phi_{x,y}(r)|\,dr \leq \kappa_1(n,M)\,d^{-n} < \infty , \tag{9}
\end{equation*}
with $\kappa_1$ depending only on $n$ and $M$.  The function $\Phi_{x,y}$ is
continuous on $(0,1)$, hence Borel there, so $(9)$ says it is Lebesgue
integrable on $(0,1)$: that is clause (1)(a).  And then
\[
  |K_{\gamma,j}(x,y)| = c_n\Bigl|\int_0^1\Phi_{x,y}\Bigr|
    \leq c_n\int_0^1|\Phi_{x,y}| \leq c_n\kappa_1\,d^{-n} ,
\]
which is clause (1)(b) with the constant $C_1 := c_n\kappa_1$.
The estimate $(9)$ on the slice integral, which is what carries clause (1)(a) and
the size bound behind (1)(b), is \noderef{GaussRieszKernelNearSliceL1Bound}; it
states exactly this bound for every $x,y$ with $|x|,|y| \leq M$ and
$0 < |x-y| \leq 2$, and is what the formalization invokes at this point.

\emph{Step 3: a uniform majorant near $x$, and clause (1)(c).}
Keep $x,y,d$ as in Step 2 and set $U = B(x,d/4)$, an open convex set.  For
$w \in U$: $|w| \leq |x| + \tfrac12 \leq R + \tfrac32 \leq M$; and
$|w-y| \geq d - d/4 = \tfrac34 d > 0$ while $|w-y| \leq \tfrac54 d \leq
\tfrac52$.  Also $w - ry = (w-y) + \tau y$ gives
$|w-ry| \geq |w-y| - \tau|y| \geq \tfrac34 d - \tau M$.  Now apply $(8)$ at the
point $w$, where its ``$d$'' is $|w-y| \leq \tfrac54 d$: the right-hand side of
$(8)$ is increasing in that quantity and decreasing in $|w-ry|$, so for all
$w \in U$ and all $\tau \in (0,1)$
\[
  |\nabla_w\Phi_{w,y}(r)| \leq H(\tau)
    := e^{M^2}\tau^{-n/2-3/2}
       \Bigl(1 + \frac{2(\tfrac54 d+\tau M)^2}{\tau}\Bigr)
       \exp\Bigl(-\frac{\bigl((\tfrac34 d - \tau M)^+\bigr)^2}{2\tau}\Bigr) ,
\]
where $s^+ = \max(s,0)$; note the right-hand side no longer depends on $w$.
Put $\tau'' = 3d/(8M)$, which satisfies $\tau'' \leq \tfrac14 < 1$ because
$d \leq 2$ and $M \geq 3$.

For $0 < \tau \leq \tau''$ we have $\tau M \leq \tfrac38 d$, hence
$\tfrac34 d - \tau M \geq \tfrac38 d$ and $\tfrac54 d + \tau M \leq 2d$, so
\[
  H(\tau) \leq e^{M^2}\Bigl(\tau^{-n/2-3/2} + 8d^2\tau^{-n/2-5/2}\Bigr)
    e^{-9d^2/(128\tau)} ,
\]
and $(4)$ with $b = 9d^2/128$, applied with $a = \tfrac{n+3}{2}$ and with
$a = \tfrac{n+5}{2}$, bounds $\int_{\{0<\tau\leq\tau''\}}H$ by
\[
  e^{M^2}\Bigl[(9/128)^{-(n+1)/2}\Gamma\bigl(\tfrac{n+1}{2}\bigr)
    + 8\,(9/128)^{-(n+3)/2}\Gamma\bigl(\tfrac{n+3}{2}\bigr)\Bigr]\,d^{-n-1} ,
\]
because $b^{1-a} = (9/128)^{-(n+1)/2}(d^2)^{-(n+1)/2}
 = (9/128)^{-(n+1)/2}d^{-n-1}$ when $a = \tfrac{n+3}{2}$, and
$8d^2\,b^{1-a} = 8\,(9/128)^{-(n+3)/2}d^2\,d^{-(n+3)} = 8\,(9/128)^{-(n+3)/2}
 d^{-n-1}$ when $a = \tfrac{n+5}{2}$.

For $\tau'' < \tau < 1$ we discard the exponential and use
$(\tfrac54 d+\tau M)^2 \leq 4d^2 + 2\tau^2M^2$, so
$H(\tau) \leq e^{M^2}\bigl(\tau^{-n/2-3/2} + 8d^2\tau^{-n/2-5/2}
 + 4M^2\tau^{-n/2-1/2}\bigr)$.  By $(5)$, with $q = \tfrac{n+3}{2}$,
$q = \tfrac{n+5}{2}$ and $q = \tfrac{n+2}{2}$ respectively (each $>1$, and in
the third case $p = \tfrac{n+1}{2} \leq q$),
\[
  \int_{\tau''}^1 H
    \leq e^{M^2}\Bigl[\tfrac{2}{n+1}(\tau'')^{-(n+1)/2}
      + \tfrac{16d^2}{n+3}(\tau'')^{-(n+3)/2}
      + \tfrac{8M^2}{n}(\tau'')^{-n/2}\Bigr] .
\]
Substituting $\tau'' = 3d/(8M)$ and using $0 < d \leq 2$ once more, the three
terms are at most constant multiples (depending only on $n$ and $M$) of
$d^{-(n+1)/2}$, $d^{2-(n+3)/2} = d^{(1-n)/2}$ and $d^{-n/2}$, and each of these
is at most $2^{(3n+2)/2}d^{-n-1}$, because
$d^{-(n+1)/2} = d^{-n-1}d^{(n+1)/2}$,
$d^{(1-n)/2} = d^{-n-1}d^{(n+3)/2}$ and
$d^{-n/2} = d^{-n-1}d^{(3n+2)/2}$, with all the displayed positive powers of
$d$ bounded by $2^{(3n+2)/2}$.  Altogether
\begin{equation*}
  \int_0^1 H(\tau)\,d\tau \leq \kappa_2(n,M)\,d^{-n-1} < \infty . \tag{10}
\end{equation*}

We now differentiate under the integral sign.  Consider
$F(w,r) = \Phi_{w,y}(r)$ on $U\times(0,1)$.  Then: (i) for each
$r \in (0,1)$, $F(\cdot,r)$ is differentiable on $U$ (Step 1); (ii) for each
$w \in U$, both $F(w,\cdot)$ and $\nabla_wF(w,\cdot)$ are measurable on
$(0,1)$, being continuous there, and $F(x,\cdot)$ is integrable by $(9)$; (iii) $|\nabla_wF(w,r)| \leq H(\tau)$ for
all $w \in U$ and $r \in (0,1)$, with $H$ integrable on $(0,1)$ by $(10)$.  By
the classical theorem on differentiation under the integral sign --- for a
convex open parameter set $U$, hypotheses (i)--(iii) imply that
$w \mapsto \int_0^1F(w,r)\,dr$ is differentiable on $U$ with derivative
$h \mapsto \int_0^1\nabla_wF(w,r)\cdot h\,dr$ --- the function
$w \mapsto \int_0^1\Phi_{w,y}(r)\,dr$ is differentiable at $x$, hence so is
$w \mapsto K_{\gamma,j}(w,y) = c_n\int_0^1\Phi_{w,y}(r)\,dr$, and its
derivative $g$ at $x$ is the linear functional
$h \mapsto c_n\int_0^1\nabla_x\Phi_{x,y}(r)\cdot h\,dr$.  Its operator norm is
$c_n\bigl|\int_0^1\nabla_x\Phi_{x,y}(r)\,dr\bigr|$, which by $(10)$ is at most
\[
  c_n\int_0^1|\nabla_x\Phi_{x,y}(r)|\,dr \leq c_n\,\kappa_2\,d^{-n-1} .
\]
This is clause (1)(c) with the constant $C_2 := c_n\kappa_2$.
The content of this step --- the uniform majorant $H$ together with the
differentiation under the integral sign that it licenses, and the resulting bound
on $\|g\|$ --- is exactly \noderef{GaussRieszKernelNearFDerivBound}, which states
this conclusion for every $x,y$ in the stated range and is what the formalization
invokes at this point.

\emph{Step 4: clause (2).}
Let $|x| \leq R$, $0 < d := |x-y| \leq 1$ and $|x-z| \leq d/2$, and let
$S = \{x + t(z-x) : t \in [0,1]\}$ be the segment from $x$ to $z$, a convex
set.  For $w \in S$ we have $|w - x| \leq |z-x| \leq d/2 \leq \tfrac12$, hence
\[
  |w| \leq |x| + \tfrac12 \leq R + 1 ,
  \qquad
  |w-y| \geq d - \tfrac{d}{2} = \tfrac{d}{2} > 0 ,
  \qquad
  |w-y| \leq d + \tfrac12 \leq \tfrac32 \leq 2 .
\]
So every $w \in S$ satisfies the hypotheses of clause (1), and clause (1)(c)
gives, for each such $w$, that $v \mapsto K_{\gamma,j}(v,y)$ is differentiable
at $w$ with derivative $g_w$ satisfying
\[
  \|g_w\| \leq \frac{C_2}{|w-y|^{\,n+1}}
    \leq \frac{C_2}{(d/2)^{\,n+1}}
    = \frac{2^{\,n+1}C_2}{d^{\,n+1}} .
\]
Differentiability at $w$ implies differentiability within $S$ at $w$ with the
same derivative.  By the mean value inequality for a function differentiable on
a convex set with derivative bounded in operator norm by a constant there
--- applied to $v \mapsto K_{\gamma,j}(v,y)$ on $S$, with the uniform bound
$2^{n+1}C_2d^{-n-1}$, and to the two points $x,z \in S$ ---
\[
  \bigl|K_{\gamma,j}(x,y) - K_{\gamma,j}(z,y)\bigr|
    \leq \frac{2^{\,n+1}C_2}{d^{\,n+1}}\,|x-z| ,
\]
which is clause (2) with the constant $2^{n+1}C_2$.

\emph{Conclusion.}  Take
$C = \max\bigl(C_1,\ C_2,\ 2^{n+1}C_2\bigr) > 0$, which depends only on $n$,
$j$ and $R$.  All the bounds proved above hold a fortiori with $C$ in place of
their individual constants, so clauses (1) and (2) hold with this single $C$.
\end{proof}
